\documentclass[12pt]{article}
\usepackage[utf8]{inputenc}
\usepackage[T1]{fontenc}
\usepackage{hyperref}
\title{Ngambay Language Resources for K-12 Curriculum Development}
\author{Compiled by Claude Sonnet 5.5 (claude-sonnet-5-5) for Andr\'as Kornai}
\date{}
\begin{document}
\maketitle

\section{Introduction}
Ngambay (Sara Ngambay, Ngambaye; ISO 639-3 \texttt{sba}; Glottocode ngam1268) is a Central Sudanic language spoken mainly in southwestern Chad around Moundou and N'Djamena, with communities in northeastern Cameroon and eastern Nigeria. Wikipedia gives about 1.38 million speakers. It is written in the Latin alphabet, with an orthography devised by the Agence de coop\'eration culturelle et technique. The language has three tones and nasalised vowels.

\section{Primary Educational Resources}
\begin{itemize}
\item ``50 Lessons in Sara-Ngambay, Volume 1'' (Thayer, Thayer, Kyamb\'e and Gondj\'e, Indiana University, 1971) is a language-learning course for outsiders. I did not see its contents. It is linked from Omniglot: \url{https://omniglot.com/writing/ngambay.htm}
\end{itemize}
I found no confirmation of Ngambay-medium school textbooks. A Chadian memoir listed in results mentions DAPLAN, the national literacy and national-languages directorate, but does not say Ngambay is a teaching language.

\section{Linguistic and Reference Materials}
\begin{itemize}
\item Ngambay dictionary by Nangone Jacob Chata, Simeon Mbayrem Djimadoum and John M. Keegan (2015), with French as the metalanguage, as listed in Glottolog: \url{https://glottolog.org/resource/reference/id/547965}
\item Palayer (1978), a fish lexicon covering Sar and Ngambay: \url{https://glottolog.org/resource/reference/id/39064}
\item Thayer (1978), a SIL study of sentence structure in Sara-Ngambay dialogues: \url{https://glottolog.org/resource/reference/id/40575}
\item Ndjerareou, Melick and Moeller (2010), ``A Brief Grammatical Sketch of Ngambay'': \url{https://grambank.clld.org/sources/sNdjerareouNgambay}
\item Grambank entry for Ngambay: \url{https://grambank.clld.org/languages/ngam1268}
\item Wikipedia overview: \url{https://en.wikipedia.org/wiki/Ngambay_language}
\end{itemize}

\section{Religious and Community Resources}
Joshua Project has a people-group page for Ngambay in Cameroon: \url{https://joshuaproject.net/index.php/people_groups/14763/CM}. Search did not confirm any Ngambay Bible translation.

\section{Digital and Audio Resources}
Wikitongues has a page for the language: \url{https://wikitongues.org/languages/sba/}. I did not verify what it holds. A bibliographic record at the CNRS REFLEX database may be relevant: \url{https://reflex.cnrs.fr/Lexiques/webball/biblio_bas.php?id=19547}

\section{Curriculum Development Context}
Ngambay is a large language with a Latin orthography and several linguistic descriptions, which is a good base. The reference description calls it a language of instruction in education, but I could not confirm this from search results. Chadian schooling is mainly in French and Arabic. The main need is verified classroom materials, which would build on the 2015 dictionary and the existing grammatical descriptions.
\end{document}